\subsection{Power and Energy Balance — Writing the System Equation}

Having established the fundamental concepts of energy storage and power dissipation, we now turn to the central task of writing complete equations that describe how physical systems behave. The power and energy balance approach provides a systematic method for deriving the differential equations governing any physical system, regardless of whether it involves electrical circuits, mechanical devices, or thermal processes. This approach rests on a single fundamental principle: energy cannot be created or destroyed, only transformed or transferred. By carefully accounting for how energy enters a system, how it is stored, and how it is dissipated, we can construct accurate mathematical models without memorizing domain-specific formulas.

The power balance equation states that the instantaneous power supplied to a system equals the rate at which energy is stored plus the rate at which energy is dissipated. In mathematical terms, if we denote the input power as $P_{\text{in}}(t)$, the rate of energy storage as $\frac{dE_{\text{stored}}}{dt}$, and the power dissipated as $P_{\text{diss}}(t)$, then the conservation equation is given by

\[
P_{\text{in}}(t) = \frac{dE_{\text{stored}}}{dt} + P_{\text{diss}}(t).
\]

This deceptively simple equation encapsulates the entire behavior of any conservative, dissipative system. The key insight is that the rate of energy storage $\frac{dE_{\text{stored}}}{dt}$ always involves the derivatives of the energy storage variables in the system. For example, in an electrical system where energy is stored in capacitors and inductors, the stored energy depends on the capacitor voltage $v_C$ and inductor current $i_L$, so its time derivative will involve $\frac{dv_C}{dt}$ and $\frac{di_L}{dt}$. Similarly, in a mechanical system, stored energy depends on position and velocity, leading to derivatives of these quantities in the balance equation.

To apply this method systematically, we follow a structured procedure that guides us from physical understanding to mathematical description. First, we identify all sources of power input to the system, which may include voltage sources, current sources, force inputs, or torque inputs depending on the domain. Second, we identify all energy storage elements and express the total stored energy as a function of the appropriate state variables. Third, we identify all dissipative elements and express the total dissipated power as a function of the appropriate system variables. Fourth, we substitute these expressions into the power balance equation and simplify to obtain a differential equation relating inputs to system states.

This systematic approach offers profound advantages over memorizing domain-specific formulas. A student who remembers that $v = L\frac{di}{dt}$ for inductors and $i = C\frac{dv}{dt}$ for capacitors might struggle when encountering a new type of system or when the circuit configuration becomes complex. However, the same student who understands that an inductor stores magnetic energy $E_L = \frac{1}{2}Li^2$ and therefore contributes a term $L i \frac{di}{dt}$ to the power balance equation can apply this same reasoning to any system involving energy storage. The power balance method derives the familiar formulas from more fundamental principles, making them easier to remember and easier to extend to novel situations.

Consider an electrical circuit containing a voltage source, a resistor, and an inductor connected in series. The voltage source provides an input voltage $v_s(t)$, the resistor has resistance $R$, and the inductor has inductance $L$. We wish to derive the differential equation relating the source voltage to the current $i(t)$ flowing through the circuit. To apply the power balance method, we first identify the power input to the system. The voltage source delivers power equal to the product of voltage and current, so $P_{\text{in}} = v_s(t) \cdot i(t)$.

Next, we identify the energy storage element. The inductor stores magnetic energy given by $E_L = \frac{1}{2}Li^2$. The rate of change of stored energy is therefore $\frac{dE_L}{dt} = \frac{d}{dt}\left(\frac{1}{2}Li^2\right) = Li \frac{di}{dt}$, where we assume $L$ is constant. The resistor dissipates power through Joule heating, with $P_{\text{diss}} = i^2R$. Substituting these expressions into the power balance equation gives

\[
v_s(t) \cdot i = Li \frac{di}{dt} + i^2R.
\]

Dividing both sides by $i$ (assuming $i \neq 0$) yields the familiar differential equation for an RL circuit:

\[
v_s(t) = L \frac{di}{dt} + Ri.
\]

This result matches the equation obtained by applying Kirchhoff's voltage law directly, but we derived it from energy principles rather than from memorized circuit rules. More importantly, the same procedure applies equally well to mechanical systems, where the analogous quantities are force, velocity, mass, and damping coefficient.

\begin{table}[htbp]
\centering
\begin{tabular}{|c|c|c|c|}
\hline
\textbf{Quantity} & \textbf{Electrical Domain} & \textbf{Mechanical (Translation)} & \textbf{Mechanical (Rotation)} \\
\hline
Input Power & $v \cdot i$ & $F \cdot v$ & $\tau \cdot \omega$ \\
\hline
Stored Energy (Inertia) & $\frac{1}{2}Li^2$ & $\frac{1}{2}mv^2$ & $\frac{1}{2}J\omega^2$ \\
\hline
Storage Rate & $Li\frac{di}{dt}$ & $mv\frac{dv}{dt}$ & $J\omega\frac{d\omega}{dt}$ \\
\hline
Dissipated Power & $i^2R$ & $Fv_{damp}$ & $\tau_{damp}\omega$ \\
\hline
\end{tabular}
\caption{Analogous quantities across energy domains, demonstrating the universal nature of power balance.}
\label{tab:domain_analogies}
\end{table}

The power balance method becomes even more powerful when applied to systems containing multiple storage elements or spanning multiple energy domains. Table~\ref{tab:domain_analogies} presents the key analogous quantities that allow us to recognize common structure across different physical domains. An electrical engineer, mechanical engineer, and aerospace engineer may use different terminology and draw different symbols on their schematics, but they are all writing equations that express the same fundamental principle of energy conservation.

Let us now work through a more complex example involving both electrical and mechanical components: an armature-controlled DC motor connected to a mechanical load. This example demonstrates how power balance naturally handles multi-domain systems without requiring us to derive conversion factors between domains. The DC motor has an armature resistance $R_a$, armature inductance $L_a$, and a back EMF constant $K_e$. The motor shaft is connected to a load with moment of inertia $J$ and viscous damping coefficient $B$. The input is an armature voltage $v_a(t)$, and the output is the rotational velocity $\omega(t)$ of the motor shaft.

The electrical power input to the armature circuit is $P_{\text{in}} = v_a(t) \cdot i_a(t)$, where $i_a(t)$ is the armature current. Energy is stored magnetically in the armature inductance, with stored energy $E_L = \frac{1}{2}L_a i_a^2$, so the rate of electrical energy storage is $L_a i_a \frac{di_a}{dt}$. Energy is also stored kinetically in the rotating inertia, with $E_J = \frac{1}{2}J\omega^2$, so the rate of mechanical energy storage is $J\omega \frac{d\omega}{dt}$.

The back EMF in the motor represents an internal conversion from electrical to mechanical energy, and it appears as a voltage drop $K_e \omega(t)$ in the electrical circuit. This means the electrical power available for storage and dissipation is actually $v_a \cdot i_a - K_e \omega \cdot i_a$, or equivalently, the input electrical power minus the rate of mechanical energy storage. Power is dissipated in the armature resistance as $i_a^2 R_a$ and in the mechanical damping as $B\omega^2$.

Writing the complete power balance requires careful attention to the energy conversion process. The net electrical power input equals the rate of electrical energy storage plus the rate of mechanical energy storage plus the total dissipated power:

\[
v_a i_a = L_a i_a \frac{di_a}{dt} + J\omega \frac{d\omega}{dt} + i_a^2 R_a + B\omega^2.
\]

The mechanical power developed by the motor, given by the torque-speed product, appears implicitly through the $K_e \omega i_a$ term representing energy conversion. The torque produced by the motor is related to the armature current by $\tau = K_t i_a$, where $K_t$ is the torque constant. For an ideal motor, $K_t = K_e$ (a property known as energy conversion reciprocity).

To derive the standard state-space model for this motor system, we must account for the relationship between electrical and mechanical power. The back EMF $K_e \omega$ represents the electrical manifestation of the mechanical motion, so we rewrite the electrical power equation as

\[
v_a i_a = L_a i_a \frac{di_a}{dt} + K_e \omega i_a + i_a^2 R_a.
\]

The term $K_e \omega i_a$ is precisely the rate at which electrical energy is converted to mechanical energy. This converted power then enters the mechanical energy balance, where it supplies the mechanical storage and dissipation:

\[
K_e \omega i_a = J\omega \frac{d\omega}{dt} + B\omega^2.
\]

Dividing both equations by the appropriate variables (assuming non-zero current and velocity) gives us the familiar coupled differential equations for the DC motor:

\[
L_a \frac{di_a}{dt} + R_a i_a + K_e \omega = v_a,
\]
\[
J \frac{d\omega}{dt} + B\omega = K_t i_a.
\]

Notice how naturally the power balance approach handles the energy conversion between domains. We did not need to memorize conversion factors or derive the motor equations from first principles of electromagnetism and mechanics; we simply accounted for energy flowing into, through, and out of the system, and the mathematics took care of itself.

The power balance method also illuminates the physical meaning of the terms in our differential equations. The $L_a \frac{di_a}{dt}$ term represents the opposition to changes in current due to energy storage in the magnetic field, while the $R_a i_a$ term represents energy dissipation as heat. The $K_e \omega$ term represents the back EMF, which is not a dissipation but rather an energy conversion mechanism. Similarly, on the mechanical side, the $J \frac{d\omega}{dt}$ term represents the rotational inertia requiring torque to accelerate, while the $B\omega$ term represents viscous damping dissipating energy as heat in the bearings and surrounding medium.

For a final worked example, let us consider a more abstract system: a mass-spring-damper system with both external forcing and a nonlinear spring. This example demonstrates that the power balance approach extends naturally to nonlinear systems and helps us understand why certain terms appear in the equations. A mass $m$ is attached to a spring with nonlinear stiffness $F_s = k_1 x + k_2 x^3$ and a damper with nonlinear friction $F_d = c_1 \dot{x} + c_2 \dot{x}^2 \operatorname{sgn}(\dot{x})$. An external force $F(t)$ acts on the mass.

The input power is $F(t) \cdot \dot{x}$. The kinetic energy stored in the moving mass is $E_K = \frac{1}{2}m\dot{x}^2$, so the rate of kinetic energy change is $m\dot{x}\ddot{x}$. The potential energy stored in the spring is $E_P = \int_0^x F_s \, dx' = \int_0^x (k_1 x' + k_2 x'^3) \, dx' = \frac{1}{2}k_1 x^2 + \frac{1}{4}k_2 x^4$, so the rate of potential energy storage is $\frac{dE_P}{dt} = (k_1 x + k_2 x^3) \dot{x}$. The power dissipated by the damper is $P_{\text{diss}} = F_d \cdot \dot{x} = (c_1 \dot{x} + c_2 \dot{x}^2 \operatorname{sgn}(\dot{x})) \dot{x} = c_1 \dot{x}^2 + c_2 |\dot{x}|^3$.

The power balance equation is therefore

\[
F \dot{x} = m\dot{x}\ddot{x} + (k_1 x + k_2 x^3)\dot{x} + c_1 \dot{x}^2 + c_2 |\dot{x}|^3.
\]

Dividing by $\dot{x}$ (for $\dot{x} \neq 0$) gives the familiar form of Newton's second law for this system:

\[
m\ddot{x} + c_1 \dot{x} + k_1 x + k_2 x^3 = F(t) - c_2 |\dot{x}|\dot{x}.
\]

The nonlinear damping term $-c_2 |\dot{x}|\dot{x}$ emerges naturally from the power balance. This term, sometimes called quadratic drag or turbulent damping, appears in many physical systems including vehicles moving through air and objects falling through viscous fluids at high Reynolds numbers. We did not need to memorize a formula for quadratic drag; it emerged directly from the energy dissipation model.

The systematic procedure for writing system equations using power balance can be summarized as follows. First, define the boundaries of your system and identify all points where energy enters or leaves the system. Second, identify all energy storage mechanisms and write expressions for the stored energy in terms of appropriate state variables. Third, identify all mechanisms of energy dissipation and write expressions for the dissipated power. Fourth, apply the power balance equation by equating the input power to the sum of storage rates and dissipation rates. Fifth, simplify the resulting equation by canceling common factors and rearranging terms to obtain a differential equation in standard form.

This procedure eliminates the need to memorize domain-specific formulas because the formulas emerge naturally from energy principles. A student who understands that inductors store magnetic energy, capacitors store electric energy, masses store kinetic energy, and springs store potential energy can derive the appropriate differential equations by simply applying the universal law of energy conservation. The power balance approach transforms differential equation derivation from a task of memorizing special cases into a task of carefully tracking how energy flows through a physical system.

As we progress through this textbook, we will encounter increasingly complex systems, including multi-degree-of-freedom mechanical systems, electrical networks with mutual coupling, thermal systems with heat transfer, and fluid systems with compressibility effects. In every case, the power balance approach provides a consistent and reliable method for deriving the governing equations. The physical intuition developed by tracking energy flows serves not only for deriving equations but also for understanding system behavior, predicting the effects of parameter changes, and designing control strategies. Energy is the common currency of physical systems, and the power balance is our accounting method for tracking it.
